\begin{center}
    \textbf{ABSTRACT}
\end{center}

\vspace{0.25cm}

Mobility analysts in Bogot\'a hold rich survey data but no direct way to turn it into microscopic traffic simulations, so every scenario requires manual, error-prone preparation of SUMO inputs. This project built an automated pipeline that converts the 2023 Mobility Survey into SUMO trips, routes and a three-dimensional road network for the Usaqu\'en locality, runs the simulation, and exposes the results in a web viewer. The viewer lets users parametrize demand (modal-shift presets, counts per mode, seed), analyze and prioritize signalized and unsignalized intersections, and compare in 3D an intersection with and without a proposed traffic light using safety indicators (TTC, DRAC, PET). The departure-time distribution of the simulation reproduces the survey profile; flow-based GEH validation remains pending.

\vspace{0.5cm}

\begin{center}
    {\textbf{RESUMEN}}
\end{center}

\vspace{0.25cm}

Los analistas de movilidad de Bogot\'a disponen de datos de encuesta, pero no de una forma directa de convertirlos en simulaciones microsc\'opicas de tr\'afico, por lo que cada escenario exige preparar a mano y con riesgo de error los insumos de SUMO. Este proyecto construy\'o un pipeline automatizado que transforma la Encuesta de Movilidad 2023 en viajes, rutas y una red vial tridimensional de la localidad de Usaqu\'en, ejecuta la simulaci\'on y publica los resultados en un visor web. El visor permite parametrizar la demanda (presets de transferencia modal, conteos por modo, semilla), analizar y priorizar intersecciones y comparar en 3D un cruce con y sin un sem\'aforo propuesto mediante indicadores de seguridad (TTC, DRAC, PET). La distribuci\'on de horas de salida simulada reproduce el perfil de la encuesta; la validaci\'on GEH con aforos queda pendiente.
